% !TeX root = ../Algebra-Kawata.tex

\section{定義・定理}

\subsection{基本的な定義}
\begin{comment}
	\begin{definition}
		% 定義を書く
	\end{definition}
	
	\begin{remark}
		% 定義に関する注意を書く
	\end{remark}
	
	\subsection{基本的な定理}
	
	\begin{theorem}
		% 定理を書く
	\end{theorem}
	
	\begin{lemma}
		% 必要なら補題を書く
	\end{lemma}
	
	\begin{corollary}
		% 必要なら系を書く
	\end{corollary}
\end{comment}

\begin{definition}[自由 $\mathbb{Z}$-加群の階数]
	$\mathbb{Z}$-加群 $M$ が階数 $n$ の自由 $\mathbb{Z}$-加群であるとは，
	ある元
	\[
	e_1,\dots,e_n\in M
	\]
	が存在して，任意の $x\in M$ が一意的に
	\[
	x=a_1e_1+\cdots+a_ne_n
	\qquad
	(a_1,\dots,a_n\in\mathbb{Z})
	\]
	と表されることをいう。
	
	このとき，$e_1,\dots,e_n$ を $M$ の $\mathbb{Z}$-基底といい，
	基底に含まれる元の個数 $n$ を $M$ の階数という。
	階数は
	\[
	\operatorname{rank}_{\mathbb{Z}} M=n
	\]
	と書く。
\end{definition}